
\documentclass{article}
\usepackage{colortbl}
\usepackage{makecell}
\usepackage{multirow}
\usepackage{supertabular}

\begin{document}

\newcounter{utterance}

\centering \large Interaction Transcript for game `bbh', experiment `boolean\_expressions', episode 71 with reflexion{-}Qwen3.5{-}27B.
\vspace{24pt}

{ \footnotesize  \setcounter{utterance}{1}
\setlength{\tabcolsep}{0pt}
\begin{supertabular}{c@{$\;$}|p{.15\linewidth}@{}p{.15\linewidth}p{.15\linewidth}p{.15\linewidth}p{.15\linewidth}p{.15\linewidth}}
    \# & \multicolumn{2}{c}{Player} && \multicolumn{2}{c}{Game Master} \\
    \hline

    \theutterance \stepcounter{utterance}  
    & & & \multicolumn{4}{p{0.6\linewidth}}{
        \cellcolor[rgb]{0.9,0.9,0.9}{
            \makecell[{{p{\linewidth}}}]{
                \texttt{\tiny{[P1$\langle$GM]}}
                \texttt{Evaluate the result of a random Boolean expression.} \\
\\ 
\texttt{Q: not ( ( not not True ) ) is} \\
\texttt{A: False} \\
\\ 
\texttt{Q: True and False and not True and True is} \\
\texttt{A: False} \\
\\ 
\texttt{Q: not not ( not ( False ) ) is} \\
\texttt{A: True} \\
\\ 
\texttt{Q: not not not not not not not True is} \\
\texttt{A:} \\
            }
        }
    }
     \\ \\

    \theutterance \stepcounter{utterance}  
    & \multicolumn{4}{p{0.6\linewidth}}{
        \cellcolor[rgb]{0.561,0.737,0.561}{
            \makecell[{{p{\linewidth}}}]{
                \texttt{\tiny{[P1$\rangle$GM]}}
                \texttt{To evaluate the expression `not not not not not not not True`, we can count the number of `not` operators applied to the value `True`.} \\
\\ 
\texttt{1.  The expression contains **7** `not` operators.} \\
\texttt{2.  Applying `not` to `True` flips the value:} \\
\texttt{    *   1 `not`: `False`} \\
\texttt{    *   2 `not`s: `True`} \\
\texttt{    *   3 `not`s: `False`} \\
\texttt{    *   ...} \\
\texttt{3.  The pattern alternates: an **even** number of `not`s results in the original value (`True`), while an **odd** number of `not`s results in the opposite value (`False`).} \\
\texttt{4.  Since 7 is an odd number, the final result is the opposite of `True`.} \\
\\ 
\texttt{**False**} \\
            }
        }
    }
    & & \\ \\

    \theutterance \stepcounter{utterance}  
    & & & \multicolumn{2}{p{0.3\linewidth}}{
        \cellcolor[rgb]{0.9,0.9,0.9}{
            \makecell[{{p{\linewidth}}}]{
                \texttt{\tiny{[GM$|$GM]}}
                \texttt{Target: False} \\
            }
        }
    }
    & & \\ \\

    \theutterance \stepcounter{utterance}  
    & & & \multicolumn{2}{p{0.3\linewidth}}{
        \cellcolor[rgb]{0.9,0.9,0.9}{
            \makecell[{{p{\linewidth}}}]{
                \texttt{\tiny{[GM$|$GM]}}
                \texttt{game\_result = LOSE} \\
            }
        }
    }
    & & \\ \\

    \theutterance \stepcounter{utterance}  
    & & \multicolumn{4}{p{0.6\linewidth}}{
        \cellcolor[rgb]{0.95,0.95,0.95}{
            \makecell[{{p{\linewidth}}}]{
                \texttt{\tiny{[P1$\leftrightarrow$P]}}
                \texttt{Post{-}episode Reflexion (not sent to GM):} \\
\texttt{New plan: My previous attempt failed because I provided a detailed step{-}by{-}step explanation instead of just the final Boolean value, which violated the implicit pattern established by the previous examples (Q: ... A: {[}Value{]}). The task requires evaluating the expression and outputting only the result ("True" or "False") without any reasoning, counting, or formatting like bold text. Next time, I will count the number of `not` operators to determine the parity (odd/even), apply the logic to the base value, and output strictly the single word result (e.g., "False") as the answer.} \\
            }
        }
    }
    & \\ \\

\end{supertabular}
}

\end{document}
